% Formatka robocza IIPF/PFMB, zgodna z Szablon_TEX_IIPF_PFMB.tex
% Kompilacja: pdfLaTeX + BibTeX + pdfLaTeX x2
\documentclass[twocolumn,a4paper,polish,10pt]{article}

\usepackage[polish]{babel}
\usepackage[utf8]{inputenc}
\usepackage[T1]{fontenc}
\usepackage{secdot}
\usepackage{indentfirst}
\usepackage{graphicx}
\usepackage{color}
\usepackage{sectsty}
\usepackage[usenames,dvipsnames,svgnames,table]{xcolor}
\usepackage[a4paper,hmargin=2cm,tmargin=1.5cm,bmargin=3cm,footskip=2cm]{geometry}
\usepackage{fancyhdr}
\usepackage{amsmath,amssymb}
\usepackage{booktabs}
\usepackage{siunitx}
\usepackage[hidelinks]{hyperref}

\sisetup{output-decimal-marker={,},group-separator={\,},detect-weight=true}
\pagestyle{fancy}
\lhead{}\chead{}\rhead{}\cfoot{}
\rfoot{\emph{K.~Nepravskaya}}
\renewcommand{\headrulewidth}{0pt}
\renewcommand{\thefootnote}{\fnsymbol{footnote}}
\cfoot{\thepage{}}
\sectionfont{\centering\large}
\graphicspath{{../figures/}}

\begin{document}

\twocolumn[
\begin{@twocolumnfalse}
\begin{center}
\textbf{\large{Tytuł ćwiczenia}}
\bigskip
\\K.~Nepravskaya\footnotemark\\
\small{opiekun: imię i nazwisko prowadzącego}\\
\textit{II Pracownia Fizyczna}\\
\small{data rozpoczęcia ćwiczenia: DD.MM.RRRR}
\end{center}

\begin{quote}\small
\begin{quote}
\begin{center}
\textbf{Streszczenie}
\end{center}
% Maksymalnie 200 słów. Podaj cel, metodę, główne wyniki liczbowe
% z niepewnościami i ewentualnie jeden krótki wniosek.
\end{quote}
\end{quote}
\bigskip
\bigskip
\end{@twocolumnfalse}
]
\footnotetext{adres e-mail}

\section{Wprowadzenie}
% Zwięzła teoria, metoda i definicje konieczne do analizy. Cytuj źródła.

\section{Opis eksperymentu}
% Aparatura, próbka, procedura i schemat układu.

\section{Wyniki i ich dyskusja}
% Obliczenia, niepewności, wykresy, tabele i interpretacja.

\section{Podsumowanie i wnioski}

\begin{thebibliography}{9}
\bibitem{instrukcja}
Instrukcja do ćwiczenia laboratoryjnego.
\end{thebibliography}

\appendix
\section{Rachunek niepewności pomiarowych}
% Szczegóły wyznaczania niepewności przenieś tutaj.

\end{document}
